\documentclass[../cuaderno6.tex]{subfiles}

\begin{document}

\subsection{Prisma}
En el caso del prisma, se puede notar una dependencia clara con el ángulo para el campo coercitivo $H_c$, visible en la figura \ref{fig:H_c_segun_H_a_Prisma_X_Fa0_03_09}. Aunque la dependencia es evidente, su forma no es clara, al contrario de lo que ocurre para la magnetización remanente $M_r$ y máxima $M_{max}$, visibles en las figuras \ref{fig:M_r_segun_H_a_Prisma_X_Fa0_03_09} y \ref{fig:M_max_segun_H_a_Prisma_X_Fa0_03_09} respectivamente, donde la magnetización decrece en cuanto se aumenta el ángulo.
\begin{figure}[H]
    \centering
    \includegraphics[width=0.5\textwidth]{../img/Fa0/Prisma X/H_c_segun_H_a_Prisma X_Fa0_03_09.png}
    \caption{El campo coercitivo $H_c$ en función del logaritmo del campo aplicado $H_{max}$ para cada ángulo.}
    \label{fig:H_c_segun_H_a_Prisma_X_Fa0_03_09}
\end{figure}
\begin{figure}[h]
    \centering
    \includegraphics[width=0.5\textwidth]{../img/Fa0/Prisma X/M_r_segun_H_a_Prisma X_Fa0_03_09.png}
    \caption{Magnetización remanente $M_r$ en función del campo aplicado $H_{max}$.}
    \label{fig:M_r_segun_H_a_Prisma_X_Fa0_03_09}
\end{figure}
\begin{figure}[h]
    \centering
    \includegraphics[width=0.5\textwidth]{../img/Fa0/Prisma X/M_max_segun_H_a_Prisma X_Fa0_03_09.png}
    \caption{Magnetización máxima $M_{max}$ en función del campo aplicado $H_{max}$.}
    \label{fig:M_max_segun_H_a_Prisma_X_Fa0_03_09}
\end{figure}

El campo coercitivo encontrado como resultado de un ajuste $H_c^0$, visible en la figura \ref{fig:H0_c_angulo_minors_Prisma}, también tiene su máximo en $75^\circ$ como el $H_c$, pero tiene un valor más bajo de lo esperado para el caso de $90^\circ$. De otra forma crecería con el ángulo. El factor de desmagnetización $N_d$, visible en la figura \ref{fig:N_d_angulo_major_Prisma_X} crece con el ángulo como es esperado, ya que al rotar se mueve a la pieza a una posición geométricamente menos conveniente para que cierren las líneas de campo sobre sus lados.

\begin{figure}[H]
    \centering
    \includegraphics[width=0.5\textwidth]{../img/Fa0/Prisma X/H0_c_angulo_minors_Prisma X.png}
    \caption{Gráfico de la pendiente del ajuste $H_c^0$ en función del ángulo.}
    \label{fig:H0_c_angulo_minors_Prisma}
\end{figure}
\begin{figure}[h]
    \centering
    \includegraphics[width=0.5\textwidth]{../img/Fa0/Prisma X/N_d_angulo_major_Prisma X.png}
    \caption{Factor de desmagnetización $N_d$ en función del ángulo.}
    \label{fig:N_d_angulo_major_Prisma_X}
\end{figure}

\end{document}